% Lösungen Einheit 2
\einheitenkopf[][2 Extrempunkte]{Einheit 2 von 5 · Extrempunkte}

\erg{15}{a) gesucht \quad b) gegeben \quad c) gegeben \quad d) gesucht \quad e) gegeben}

\erg{16}{a) (1) \quad b) (2), denn $f''$ entscheidet hier nicht \quad c) (4) \quad d) (3)}

\erg{17}{a) $x = -3$, $x = 3$ \quad b) $f'(x) = 3(x + 3)(x - 1)$: $x = -3$, $x = 1$ \quad c) $f'(x) = 6(x - 1)(x - 2)$: $x = 1$, $x = 2$ \quad d) $f'(x) = -3(x - 1)(x - 3)$: $x = 1$, $x = 3$ \quad e) $f'(x) = 3(x + 1)(x - 3)$, $f''(x) = 6x - 6$: $H(-1\,|\,7)$, $T(3\,|\,{-}25)$ \quad f) $f'(x) = 12x\,(x + 1)(x - 2)$, $f''(x) = 36x^2 - 24x - 24$: $T(-1\,|\,{-}5)$, $H(0\,|\,0)$, $T(2\,|\,{-}32)$ \quad g) $f'(x) = 12x^2(x + 1)$: $T(-1\,|\,{-}1)$; bei $x = 0$ ist $f''(0) = 0$ und $f'$ wechselt das Vorzeichen nicht (oder $f'''(0) = 24 \ne 0$): Sattelpunkt $S(0\,|\,0)$ \quad h) $f'(x) = 4x^3 - 12x^2 - 4x + 12 = 4(x - 3)(x^2 - 1)$, $f''(x) = 12x^2 - 24x - 4$: $T(-1\,|\,{-}9)$, $H(1\,|\,7)$, $T(3\,|\,{-}9)$}

\erg{18}{a) $f'(x) = (1 - x)\,e^{-x}$, $f''(x) = (x - 2)\,e^{-x}$, $f''(1) < 0$: $H(1\,|\,e^{-1})$ \quad b) $k'(t) = 8e^{-0{,}5t}(1 - 0{,}5t) = 0$ für $t = 2$; $k(2) = 16e^{-1} \approx 5{,}89$; nach 2 Stunden ist die Konzentration mit etwa 5,89 mg pro Liter am höchsten – Hochpunkt, weil die Konzentration bei $t = 0$ null ist, danach steigt und langfristig wieder gegen null geht, und es nur diese eine Stelle mit $k' = 0$ gibt \quad c) $H(-1\,|\,7)$, $T(2\,|\,{-}20)$, Randwerte $f(-2) = -4$, $f(4) = 32$: größter Wert 32 am rechten Rand \quad d) $f'(x) = 4x^3 - 32 = 0$ für $x = 2$, $T(2\,|\,{-}48)$; für $x \to \pm\infty$ geht $f(x) \to +\infty$: Wertebereich $[-48;\,\infty[$ \quad e) $f'(x) = (-x^2 + 2x + 3)\,e^{-x} = -(x + 1)(x - 3)\,e^{-x}$, $f''(x) = (x^2 - 4x - 1)\,e^{-x}$; $f''(-1) = 4e > 0$: $T(-1\,|\,{-}2e)$; $f''(3) = -4e^{-3} < 0$: $H(3\,|\,6e^{-3}) \approx H(3\,|\,0{,}30)$}

\erg{19}{a) $f'(4) = 3\cdot 16 - 48 = 0$ \quad b) $f'(5) = 10 - 10 = 0$ \quad c) $f'(1) = 4 - 4 = 0$ \quad d) $f'(-1) = -6 + 6 = 0$ \quad e) $f'(-2) = 0$, $f''(-2) = -6 < 0$: Hochpunkt $(-2\,|\,2)$ \quad f) $f'(0) = 0$, $f''(0) = 0$; $f'(x) = 4x^3$ wechselt bei 0 von $-$ nach $+$: Tiefpunkt $(0\,|\,{-}5)$ \quad g) $f'(2) = -12 + 12 = 0$, $f''(2) = -12 < 0$, $f(2) = -8 + 24 + 3 = 19$ \quad h) $f''(1) < 0$ und $f'(1) = 0$: Hochpunkt bei $x = 1$ \quad i) $f'(1) = 0$, $f''(1) = 0$, $f'(x) = 3(x - 1)^2 \ge 0$ wechselt das Vorzeichen nicht: kein Extrempunkt, Sattelpunkt $S(1\,|\,1)$ \quad j) $f'$ wechselt zwischen 1 und 2 von negativ nach positiv; an der einzigen Nullstelle dazwischen fällt $f$ vorher und steigt danach: Tiefstelle \quad k) $f'(x) = -3x^2 + 6x$, $f'(2) = 0$; $f''(x) = -6x + 6$, $f''(2) = -6 < 0$; $f(2) = -8 + 12 + 1 = 5$}

\erg{20}{a) Sara deutet das Vorzeichen von $f''$ falsch: $f''(1) > 0$ heißt Tiefpunkt. Richtig: $T(1\,|\,{-}3)$ und $H(-1\,|\,5)$ \quad b) Ben setzt $e^{x} = 0$, aber $e^{x}$ ist nie null. Richtig: nur $x = 1$; $f''(x) = x\,e^{x}$, $f''(1) = e > 0$: $T(1\,|\,{-}e)$}

\erg{21}{a) Nein: $f'(x_0) = 0$ ist nur notwendig; an einer Sattelstelle wie bei $f(x) = x^3 + 1$ in $0$ ist $f'$ null, aber kein Extrempunkt. \quad b) Ja: Der größte Wert kann am Rand liegen, dort muss die Tangente nicht waagerecht sein (z.\,B. wenn $f$ auf $[0;\,4]$ steigt, bei $x = 4$).}

\erg{22}{a) $h'(x) = -0{,}2(x - 1)(x - 4)(x - 7)$; Hochpunkte $H_1(1\,|\,4{,}45)$ und $H_2(7\,|\,4{,}45)$, dazwischen $T(4\,|\,0{,}4)$; also Kuppen 44,5 m hoch bei 10 m und 70 m \quad b) Abstand $7 - 1 = 6$ Längeneinheiten, das sind 60 m (gleiche Höhe) – die Bedingung ist erfüllt.}
